This section identifies the most promising directions for extending the Codynamic Book Machine beyond the current implementation. The present system already demonstrates durable work objects, explicit task queues, specialist agents, and a reviewable artifact pipeline, but several capabilities remain intentionally underdeveloped. The most immediate opportunities are to make diagram planning more semantically aware, to improve how revision work is scheduled across the whole manuscript, to strengthen retrieval and registration of references, to expose more of the runtime state in the user interface, and to evaluate the system more formally against repeatable criteria.

First, diagram support could move from opportunistic requests toward richer planning. At present, section agents can propose one or two diagrams when a visual would materially clarify a section, but the decision is local and largely manual. A more capable planner could inspect the outline, the dependency graph, and prior diagram memory to identify recurring structures that deserve consistent visualization across chapters. For example, the system could distinguish between architecture diagrams, workflow diagrams, and state-transition diagrams, and it could avoid redundant figures by comparing proposed diagrams against previously generated assets. Such a planner would make diagram generation less reactive and more aligned with the manuscript's conceptual structure.

Second, revision scheduling could become more global and less purely section-local. The current gardening cycle is effective at catching local drift, repeated claims, and unsupported assertions, but it still depends on targeted callbacks and incremental maintenance. A future scheduler could rank revision work across the entire book by severity, dependency depth, and downstream impact, then batch related edits so that cross-references, terminology, and narrative transitions are repaired together. This would be especially useful when a change in one chapter affects definitions or phrasing in later chapters. A global scheduler could also reduce churn by preventing multiple agents from revising the same conceptual area in isolation.

Third, reference retrieval and registration could be made more robust. The present workflow correctly separates citation ownership from section drafting, but it still relies on the agent to notice missing support and to request research at the right time. Future work could add stronger claim-to-source matching, better detection of unsupported technical statements, and more structured handoff records for references\_agent. In addition, the system could surface citation gaps earlier in the drafting process, before prose is finalized, so that unsupported claims are less likely to survive into later review stages. This would improve both factual reliability and the efficiency of the editorial loop.

Fourth, the user interface could expose more of the system's internal state without overwhelming the reader. Because the machine is built around durable queues, append-only messages, and generated artifacts, it is well suited to observability. A richer UI could show task status, pending callbacks, revision history, diagram requests, citation registration state, and build outcomes in a single place. It could also make it easier to inspect why a section changed, which agent made the change, and what evidence or feedback triggered the revision. Better observability would not only help operators debug the system; it would also make the architecture easier to trust and explain.

Finally, the system would benefit from more formal evaluation. The current paper is grounded in a working repository, but the evidence is still primarily qualitative and case-study based. Future work could define benchmarks for revision quality, citation completeness, diagram usefulness, build stability, and turnaround time across repeated editing cycles. It could also compare the Codynamic Book Machine against simpler pipelines that lack durable queues or specialist agents. Such evaluation would help distinguish which parts of the architecture are essential, which are merely convenient, and which produce measurable gains in long-form authoring.

Taken together, these directions point toward a more autonomous, more legible, and more empirically grounded book machine. The current system already shows that books can be treated as coordinated technical artifacts; future work is about making that coordination deeper, more observable, and easier to validate.
